\chapter{Tableaux et structures de données}

\textit{Les tableaux sont des structures de données fondamentales en programmation. Ce chapitre présente les tableaux à une et deux dimensions, ainsi que les chaînes de caractères, avec leurs algorithmes de base : parcours, recherche, tri et comptage.}

\section{L'essentiel du cours}

\subsection{Tableaux à une dimension}

\begin{definition}[Tableau]
Un tableau est une structure de données qui permet de stocker plusieurs valeurs de même type dans une variable unique. Chaque élément est repéré par un indice (ou index). En algorithmique, on déclare un tableau par : \texttt{Tableau nom[taille] : Type}.
\end{definition}

\begin{propriete}
\begin{itemize}
    \item Le premier élément a l'indice 1 (ou 0 selon le langage).
    \item La taille est fixée à la déclaration.
    \item On accède à un élément par \texttt{nom[indice]}.
\end{itemize}
\end{propriete}

\begin{aretenir}
\begin{itemize}
    \item \textbf{Parcours complet} : Pour i de 1 à N faire ... FinPour
    \item \textbf{Recherche d'un élément} : parcourir jusqu'à trouver la valeur.
    \item \textbf{Maximum/Minimum} : initialiser avec le premier élément, puis comparer.
    \item \textbf{Somme/Moyenne} : accumuler dans une variable.
    \item \textbf{Tri simple} : tri par sélection ou par insertion.
\end{itemize}
\end{aretenir}

\begin{exemple}[Parcours et somme]
\begin{algo}
Algorithme SommeTableau
Variables
    T : Tableau[100] : Réel
    N, i : Entier
    S : Réel
Début
    Ecrire("Donner la taille N :")
    Lire(N)
    Pour i de 1 à N faire
        Ecrire("T[", i, "] = ")
        Lire(T[i])
    FinPour
    S ← 0
    Pour i de 1 à N faire
        S ← S + T[i]
    FinPour
    Ecrire("Somme = ", S)
Fin
\end{algo}
\end{exemple}

\subsection{Tableaux à deux dimensions}

\begin{definition}[Matrice]
Un tableau à deux dimensions (ou matrice) se déclare par : \texttt{Tableau nom[N][M] : Type}. On accède à un élément par \texttt{nom[i][j]}, où i est la ligne et j la colonne.
\end{definition}

\begin{aretenir}
Pour parcourir une matrice, on utilise deux boucles imbriquées :
\begin{itemize}
    \item Boucle externe : parcours des lignes.
    \item Boucle interne : parcours des colonnes.
\end{itemize}
\end{aretenir}

\subsection{Chaînes de caractères}

\begin{definition}[Chaîne]
Une chaîne de caractères est une suite de caractères. En algorithmique, on la manipule comme un tableau de caractères. La longueur d'une chaîne est donnée par \texttt{Longueur(ch)}.
\end{definition}

\begin{propriete}
Opérations courantes sur les chaînes :
\begin{itemize}
    \item \texttt{Concaténer(ch1, ch2)} : fusionner deux chaînes.
    \item \texttt{SousChaîne(ch, début, longueur)} : extraire une partie.
    \item \texttt{PosCar(ch, car)} : position d'un caractère.
    \item \texttt{Effacer(ch, pos, nb)} : supprimer des caractères.
\end{itemize}
\end{propriete}

\section{Exercices}

\rubrique{Tableaux à une dimension : parcours et opérations de base}

\exo{}\diff{1}
Écrire un algorithme qui remplit un tableau de 10 entiers saisis au clavier, puis affiche le nombre d'éléments pairs.

\exo{}\diff{1}
Écrire un algorithme qui calcule la moyenne des notes d'une classe de N élèves (N ≤ 50). Les notes sont stockées dans un tableau.

\exo{}\diff{1}
Écrire un algorithme qui recherche la valeur maximale dans un tableau de N entiers.

\exo{}\diff{1}
Écrire un algorithme qui recherche la valeur minimale dans un tableau de N réels.

\exo{}\diff{1}
Écrire un algorithme qui calcule la somme des éléments d'un tableau de N entiers.

\exo{}\diff{2}
Écrire un algorithme qui lit un tableau de N entiers, puis affiche le nombre d'occurrences d'une valeur donnée V.

\exo{}\diff{2}
Écrire un algorithme qui vérifie si un tableau de N entiers est trié par ordre croissant.

\exo{}\diff{2}
Écrire un algorithme qui inverse l'ordre des éléments d'un tableau de N entiers.

\exo{}\diff{2}
Écrire un algorithme qui supprime toutes les occurrences d'une valeur V dans un tableau de N entiers.

\exo{}\diff{2}
Écrire un algorithme qui insère une valeur V à une position P dans un tableau de N entiers.

\rubrique{Tri et recherche}

\exo{}\diff{2}
Écrire un algorithme de tri par sélection pour trier un tableau de N entiers par ordre croissant.

\exo{}\diff{2}
Écrire un algorithme de tri par insertion pour trier un tableau de N entiers par ordre décroissant.

\exo{}\diff{2}
Écrire un algorithme de recherche dichotomique dans un tableau trié de N entiers.

\exo{}\diff{3}
On dispose d'un tableau T de N entiers. Écrire un algorithme qui affiche les éléments uniques (qui apparaissent une seule fois).

\exo{}\diff{3}
Écrire un algorithme qui fusionne deux tableaux triés T1 et T2 en un seul tableau trié T3.

\rubrique{Tableaux à deux dimensions}

\exo{}\diff{1}
Écrire un algorithme qui remplit une matrice M de taille N×M avec des valeurs saisies au clavier, puis affiche la matrice.

\exo{}\diff{2}
Écrire un algorithme qui calcule la somme des éléments d'une matrice N×M.

\exo{}\diff{2}
Écrire un algorithme qui affiche les éléments de la diagonale principale d'une matrice carrée N×N.

\exo{}\diff{2}
Écrire un algorithme qui recherche le maximum dans une matrice N×M.

\exo{}\diff{3}
Écrire un algorithme qui calcule le produit de deux matrices A (N×P) et B (P×M).

\rubrique{Chaînes de caractères}

\exo{}\diff{1}
Écrire un algorithme qui lit une chaîne de caractères et affiche sa longueur.

\exo{}\diff{2}
Écrire un algorithme qui lit une chaîne et affiche le nombre de voyelles qu'elle contient.

\exo{}\diff{2}
Écrire un algorithme qui vérifie si une chaîne est un palindrome (se lit identique dans les deux sens).

\exo{}\diff{2}
Écrire un algorithme qui supprime les espaces d'une chaîne de caractères.

\exo{}\diff{3}
Écrire un algorithme qui lit une phrase et affiche chaque mot sur une ligne séparée.

\rubrique{Exercices type Baccalauréat}

\sujetbac[2023, Cameroun, C]
On considère un tableau T de N entiers (N ≤ 100). Écrire un algorithme qui :
\begin{enumerate}
    \item Remplit le tableau avec des valeurs saisies au clavier.
    \item Calcule et affiche la somme des éléments positifs.
    \item Calcule et affiche le produit des éléments négatifs.
    \item Affiche le nombre d'éléments nuls.
\end{enumerate}

\sujetbac[2022, Cameroun, D]
Un professeur souhaite gérer les notes de ses 30 élèves. Écrire un algorithme qui :
\begin{enumerate}
    \item Saisit les notes dans un tableau.
    \item Calcule la moyenne de la classe.
    \item Affiche le nombre d'élèves ayant une note supérieure ou égale à 10.
    \item Affiche la note maximale et la note minimale.
\end{enumerate}

\sujetbac[2021, Cameroun, E]
On dispose d'un tableau T de N entiers. Écrire un algorithme qui :
\begin{enumerate}
    \item Trie le tableau par ordre croissant en utilisant le tri par sélection.
    \item Demande une valeur V à rechercher.
    \item Affiche la position de V dans le tableau trié (ou un message si absent).
\end{enumerate}

\section{Corrigés}

\corrige{de l'exercice 1}
\begin{algo}
Algorithme CompterPairs
Variables
    T : Tableau[10] : Entier
    i, nbPairs : Entier
Début
    nbPairs ← 0
    Pour i de 1 à 10 faire
        Ecrire("T[", i, "] = ")
        Lire(T[i])
        Si T[i] mod 2 = 0 Alors
            nbPairs ← nbPairs + 1
        FinSi
    FinPour
    Ecrire("Nombre d'éléments pairs : ", nbPairs)
Fin
\end{algo}

\corrige{de l'exercice 2}
\begin{algo}
Algorithme MoyenneClasse
Variables
    notes : Tableau[50] : Réel
    N, i : Entier
    somme, moyenne : Réel
Début
    Ecrire("Nombre d'élèves : ")
    Lire(N)
    somme ← 0
    Pour i de 1 à N faire
        Ecrire("Note de l'élève ", i, " : ")
        Lire(notes[i])
        somme ← somme + notes[i]
    FinPour
    moyenne ← somme / N
    Ecrire("Moyenne de la classe : ", moyenne)
Fin
\end{algo}

\corrige{de l'exercice 3}
\begin{algo}
Algorithme MaximumTableau
Variables
    T : Tableau[100] : Entier
    N, i, max : Entier
Début
    Ecrire("Taille du tableau : ")
    Lire(N)
    Pour i de 1 à N faire
        Ecrire("T[", i, "] = ")
        Lire(T[i])
    FinPour
    max ← T[1]
    Pour i de 2 à N faire
        Si T[i] > max Alors
            max ← T[i]
        FinSi
    FinPour
    Ecrire("Maximum : ", max)
Fin
\end{algo}

\corrige{de l'exercice 4}
\begin{algo}
Algorithme MinimumTableau
Variables
    T : Tableau[100] : Réel
    N, i : Entier
    min : Réel
Début
    Ecrire("Taille du tableau : ")
    Lire(N)
    Pour i de 1 à N faire
        Ecrire("T[", i, "] = ")
        Lire(T[i])
    FinPour
    min ← T[1]
    Pour i de 2 à N faire
        Si T[i] < min Alors
            min ← T[i]
        FinSi
    FinPour
    Ecrire("Minimum : ", min)
Fin
\end{algo}

\corrige{de l'exercice 5}
\begin{algo}
Algorithme SommeTableau
Variables
    T : Tableau[100] : Entier
    N, i, somme : Entier
Début
    Ecrire("Taille du tableau : ")
    Lire(N)
    somme ← 0
    Pour i de 1 à N faire
        Ecrire("T[", i, "] = ")
        Lire(T[i])
        somme ← somme + T[i]
    FinPour
    Ecrire("Somme : ", somme)
Fin
\end{algo}

\corrige{de l'exercice 6}
\begin{algo}
Algorithme Occurrences
Variables
    T : Tableau[100] : Entier
    N, i, V, compteur : Entier
Début
    Ecrire("Taille du tableau : ")
    Lire(N)
    Pour i de 1 à N faire
        Ecrire("T[", i, "] = ")
        Lire(T[i])
    FinPour
    Ecrire("Valeur à rechercher : ")
    Lire(V)
    compteur ← 0
    Pour i de 1 à N faire
        Si T[i] = V Alors
            compteur ← compteur + 1
        FinSi
    FinPour
    Ecrire("Nombre d'occurrences de ", V, " : ", compteur)
Fin
\end{algo}

\corrige{de l'exercice 7}
\begin{algo}
Algorithme VerifierTri
Variables
    T : Tableau[100] : Entier
    N, i : Entier
    trie : Booléen
Début
    Ecrire("Taille du tableau : ")
    Lire(N)
    Pour i de 1 à N faire
        Ecrire("T[", i, "] = ")
        Lire(T[i])
    FinPour
    trie ← Vrai
    Pour i de 1 à N-1 faire
        Si T[i] > T[i+1] Alors
            trie ← Faux
        FinSi
    FinPour
    Si trie Alors
        Ecrire("Le tableau est trié par ordre croissant.")
    Sinon
        Ecrire("Le tableau n'est pas trié.")
    FinSi
Fin
\end{algo}

\corrige{de l'exercice 8}
\begin{algo}
Algorithme InverserTableau
Variables
    T : Tableau[100] : Entier
    N, i, temp : Entier
Début
    Ecrire("Taille du tableau : ")
    Lire(N)
    Pour i de 1 à N faire
        Ecrire("T[", i, "] = ")
        Lire(T[i])
    FinPour
    Pour i de 1 à N div 2 faire
        temp ← T[i]
        T[i] ← T[N - i + 1]
        T[N - i + 1] ← temp
    FinPour
    Ecrire("Tableau inversé :")
    Pour i de 1 à N faire
        Ecrire(T[i])
    FinPour
Fin
\end{algo}

\corrige{de l'exercice 9}
\begin{algo}
Algorithme SupprimerValeur
Variables
    T : Tableau[100] : Entier
    N, i, j, V : Entier
Début
    Ecrire("Taille du tableau : ")
    Lire(N)
    Pour i de 1 à N faire
        Ecrire("T[", i, "] = ")
        Lire(T[i])
    FinPour
    Ecrire("Valeur à supprimer : ")
    Lire(V)
    i ← 1
    TantQue i ≤ N faire
        Si T[i] = V Alors
            Pour j de i à N-1 faire
                T[j] ← T[j+1]
            FinPour
            N ← N - 1
        Sinon
            i ← i + 1
        FinSi
    FinTantQue
    Ecrire("Nouveau tableau :")
    Pour i de 1 à N faire
        Ecrire(T[i])
    FinPour
Fin
\end{algo}

\corrige{de l'exercice 10}
\begin{algo}
Algorithme InsererValeur
Variables
    T : Tableau[101] : Entier
    N, i, V, P : Entier
Début
    Ecrire("Taille du tableau : ")
    Lire(N)
    Pour i de 1 à N faire
        Ecrire("T[", i, "] = ")
        Lire(T[i])
    FinPour
    Ecrire("Valeur à insérer : ")
    Lire(V)
    Ecrire("Position d'insertion : ")
    Lire(P)
    Pour i de N jusqu'à P pas -1 faire
        T[i+1] ← T[i]
    FinPour
    T[P] ← V
    N ← N + 1
    Ecrire("Nouveau tableau :")
    Pour i de 1 à N faire
        Ecrire(T[i])
    FinPour
Fin
\end{algo}

\corrige{de l'exercice 11}
\begin{algo}
Algorithme TriSelection
Variables
    T : Tableau[100] : Entier
    N, i, j, min, temp : Entier
Début
    Ecrire("Taille du tableau : ")
    Lire(N)
    Pour i de 1 à N faire
        Ecrire("T[", i, "] = ")
        Lire(T[i])
    FinPour
    Pour i de 1 à N-1 faire
        min ← i
        Pour j de i+1 à N faire
            Si T[j] < T[min] Alors
                min ← j
            FinSi
        FinPour
        temp ← T[i]
        T[i] ← T[min]
        T[min] ← temp
    FinPour
    Ecrire("Tableau trié :")
    Pour i de 1 à N faire
        Ecrire(T[i])
    FinPour
Fin
\end{algo}

\corrige{de l'exercice 12}
\begin{algo}
Algorithme TriInsertion
Variables
    T : Tableau[100] : Entier
    N, i, j, cle : Entier
Début
    Ecrire("Taille du tableau : ")
    Lire(N)
    Pour i de 1 à N faire
        Ecrire("T[", i, "] = ")
        Lire(T[i])
    FinPour
    Pour i de 2 à N faire
        cle ← T[i]
        j ← i - 1
        TantQue j ≥ 1 et T[j] < cle faire
            T[j+1] ← T[j]
            j ← j - 1
        FinTantQue
        T[j+1] ← cle
    FinPour
    Ecrire("Tableau trié par ordre décroissant :")
    Pour i de 1 à N faire
        Ecrire(T[i])
    FinPour
Fin
\end{algo}

\corrige{de l'exercice 13}
\begin{algo}
Algorithme RechercheDichotomique
Variables
    T : Tableau[100] : Entier
    N, i, V, debut, fin, milieu : Entier
    trouve : Booléen
Début
    Ecrire("Taille du tableau : ")
    Lire(N)
    Pour i de 1 à N faire
        Ecrire("T[", i, "] = ")
        Lire(T[i])
    FinPour
    Ecrire("Valeur à rechercher : ")
    Lire(V)
    debut ← 1
    fin ← N
    trouve ← Faux
    TantQue debut ≤ fin et non trouve faire
        milieu ← (debut + fin) div 2
        Si T[milieu] = V Alors
            trouve ← Vrai
        Sinon
            Si T[milieu] > V Alors
                fin ← milieu - 1
            Sinon
                debut ← milieu + 1
            FinSi
        FinSi
    FinTantQue
    Si trouve Alors
        Ecrire("Valeur trouvée à la position ", milieu)
    Sinon
        Ecrire("Valeur non trouvée")
    FinSi
Fin
\end{algo}

\corrige{de l'exercice 14}
\begin{algo}
Algorithme ElementsUniques
Variables
    T : Tableau[100] : Entier
    N, i, j : Entier
    estUnique : Booléen
Début
    Ecrire("Taille du tableau : ")
    Lire(N)
    Pour i de 1 à N faire
        Ecrire("T[", i, "] = ")
        Lire(T[i])
    FinPour
    Ecrire("Éléments uniques :")
    Pour i de 1 à N faire
        estUnique ← Vrai
        Pour j de 1 à N faire
            Si i ≠ j et T[i] = T[j] Alors
                estUnique ← Faux
            FinSi
        FinPour
        Si estUnique Alors
            Ecrire(T[i])
        FinSi
    FinPour
Fin
\end{algo}

\corrige{de l'exercice 15}
\begin{algo}
Algorithme FusionTableaux
Variables
    T1, T2, T3 : Tableau[100] : Entier
    N1, N2, i, j, k : Entier
Début
    Ecrire("Taille du premier tableau : ")
    Lire(N1)
    Pour i de 1 à N1 faire
        Ecrire("T1[", i, "] = ")
        Lire(T1[i])
    FinPour
    Ecrire("Taille du deuxième tableau : ")
    Lire(N2)
    Pour i de 1 à N2 faire
        Ecrire("T2[", i, "] = ")
        Lire(T2[i])
    FinPour
    i ← 1
    j ← 1
    k ← 1
    TantQue i ≤ N1 et j ≤ N2 faire
        Si T1[i] < T2[j] Alors
            T3[k] ← T1[i]
            i ← i + 1
        Sinon
            T3[k] ← T2[j]
            j ← j + 1
        FinSi
        k ← k + 1
    FinTantQue
    TantQue i ≤ N1 faire
        T3[k] ← T1[i]
        i ← i + 1
        k ← k + 1
    FinTantQue
    TantQue j ≤ N2 faire
        T3[k] ← T2[j]
        j ← j + 1
        k ← k + 1
    FinTantQue
    Ecrire("Tableau fusionné trié :")
    Pour i de 1 à k-1 faire
        Ecrire(T3[i])
    FinPour
Fin
\end{algo}

\corrige{de l'exercice 16}
\begin{algo}
Algorithme AfficherMatrice
Variables
    M : Tableau[50][50] : Réel
    N, Mdim, i, j : Entier
Début
    Ecrire("Nombre de lignes : ")
    Lire(N)
    Ecrire("Nombre de colonnes : ")
    Lire(Mdim)
    Pour i de 1 à N faire
        Pour j de 1 à Mdim faire
            Ecrire("M[", i, "][", j, "] = ")
            Lire(M[i][j])
        FinPour
    FinPour
    Ecrire("Matrice :")
    Pour i de 1 à N faire
        Pour j de 1 à Mdim faire
            Ecrire(M[i][j])
        FinPour
        Ecrire("")  {saut de ligne}
    FinPour
Fin
\end{algo}

\corrige{de l'exercice 17}
\begin{algo}
Algorithme SommeMatrice
Variables
    M : Tableau[50][50] : Réel
    N, Mdim, i, j : Entier
    somme : Réel
Début
    Ecrire("Nombre de lignes : ")
    Lire(N)
    Ecrire("Nombre de colonnes : ")
    Lire(Mdim)
    Pour i de 1 à N faire
        Pour j de 1 à Mdim faire
            Ecrire("M[", i, "][", j, "] = ")
            Lire(M[i][j])
        FinPour
    FinPour
    somme ← 0
    Pour i de 1 à N faire
        Pour j de 1 à Mdim faire
            somme ← somme + M[i][j]
        FinPour
    FinPour
    Ecrire("Somme des éléments : ", somme)
Fin
\end{algo}

\corrige{de l'exercice 18}
\begin{algo}
Algorithme DiagonalePrincipale
Variables
    M : Tableau[50][50] : Réel
    N, i : Entier
Début
    Ecrire("Taille de la matrice carrée : ")
    Lire(N)
    Pour i de 1 à N faire
        Pour j de 1 à N faire
            Ecrire("M[", i, "][", j, "] = ")
            Lire(M[i][j])
        FinPour
    FinPour
    Ecrire("Diagonale principale :")
    Pour i de 1 à N faire
        Ecrire(M[i][i])
    FinPour
Fin
\end{algo}

\corrige{de l'exercice 19}
\begin{algo}
Algorithme MaximumMatrice
Variables
    M : Tableau[50][50] : Réel
    N, Mdim, i, j : Entier
    max : Réel
Début
    Ecrire("Nombre de lignes : ")
    Lire(N)
    Ecrire("Nombre de colonnes : ")
    Lire(Mdim)
    Pour i de 1 à N faire
        Pour j de 1 à Mdim faire
            Ecrire("M[", i, "][", j, "] = ")
            Lire(M[i][j])
        FinPour
    FinPour
    max ← M[1][1]
    Pour i de 1 à N faire
        Pour j de 1 à Mdim faire
            Si M[i][j] > max Alors
                max ← M[i][j]
            FinSi
        FinPour
    FinPour
    Ecrire("Maximum : ", max)
Fin
\end{algo}

\corrige{de l'exercice 20}
\begin{algo}
Algorithme ProduitMatrices
Variables
    A, B, C : Tableau[50][50] : Réel
    N, P, M, i, j, k : Entier
    somme : Réel
Début
    Ecrire("Dimensions de A (N×P) : ")
    Lire(N, P)
    Ecrire("Dimensions de B (P×M) : ")
    Lire(P, M)
    Pour i de 1 à N faire
        Pour j de 1 à P faire
            Ecrire("A[", i, "][", j, "] = ")
            Lire(A[i][j])
        FinPour
    FinPour
    Pour i de 1 à P faire
        Pour j de 1 à M faire
            Ecrire("B[", i, "][", j, "] = ")
            Lire(B[i][j])
        FinPour
    FinPour
    Pour i de 1 à N faire
        Pour j de 1 à M faire
            somme ← 0
            Pour k de 1 à P faire
                somme ← somme + A[i][k] * B[k][j]
            FinPour
            C[i][j] ← somme
        FinPour
    FinPour
    Ecrire("Matrice produit C :")
    Pour i de 1 à N faire
        Pour j de 1 à M faire
            Ecrire(C[i][j])
        FinPour
        Ecrire("")
    FinPour
Fin
\end{algo}

\corrige{de l'exercice 21}
\begin{algo}
Algorithme LongueurChaine
Variables
    ch : Chaîne
    longueur : Entier
Début
    Ecrire("Donner une chaîne : ")
    Lire(ch)
    longueur ← Longueur(ch)
    Ecrire("Longueur de la chaîne : ", longueur)
Fin
\end{algo}

\corrige{de l'exercice 22}
\begin{algo}
Algorithme CompterVoyelles
Variables
    ch : Chaîne
    i, nbVoyelles, longueur : Entier
    car : Caractère
Début
    Ecrire("Donner une chaîne : ")
    Lire(ch)
    longueur ← Longueur(ch)
    nbVoyelles ← 0
    Pour i de 1 à longueur faire
        car ← SousChaîne(ch, i, 1)
        Si car = 'a' ou car = 'e' ou car = 'i' ou car = 'o' ou car = 'u' ou
           car = 'A' ou car = 'E' ou car = 'I' ou car = 'O' ou car = 'U' Alors
            nbVoyelles ← nbVoyelles + 1
        FinSi
    FinPour
    Ecrire("Nombre de voyelles : ", nbVoyelles)
Fin
\end{algo}

\corrige{de l'exercice 23}
\begin{algo}
Algorithme Palindrome
Variables
    ch : Chaîne
    i, longueur : Entier
    estPalindrome : Booléen
Début
    Ecrire("Donner une chaîne : ")
    Lire(ch)
    longueur ← Longueur(ch)
    estPalindrome ← Vrai
    Pour i de 1 à longueur div 2 faire
        Si SousChaîne(ch, i, 1) ≠ SousChaîne(ch, longueur - i + 1, 1) Alors
            estPalindrome ← Faux
        FinSi
    FinPour
    Si estPalindrome Alors
        Ecrire("La chaîne est un palindrome.")
    Sinon
        Ecrire("La chaîne n'est pas un palindrome.")
    FinSi
Fin
\end{algo}

\corrige{de l'exercice 24}
\begin{algo}
Algorithme SupprimerEspaces
Variables
    ch, resultat : Chaîne
    i, longueur : Entier
    car : Caractère
Début
    Ecrire("Donner une chaîne : ")
    Lire(ch)
    longueur ← Longueur(ch)
    resultat ← ""
    Pour i de 1 à longueur faire
        car ← SousChaîne(ch, i, 1)
        Si car ≠ ' ' Alors
            resultat ← Concaténer(resultat, car)
        FinSi
    FinPour
    Ecrire("Chaîne sans espaces : ", resultat)
Fin
\end{algo}

\corrige{de l'exercice 25}
\begin{algo}
Algorithme DecouperMots
Variables
    phrase, mot : Chaîne
    i, longueur : Entier
    car : Caractère
Début
    Ecrire("Donner une phrase : ")
    Lire(phrase)
    longueur ← Longueur(phrase)
    mot ← ""
    Pour i de 1 à longueur faire
        car ← SousChaîne(phrase, i, 1)
        Si car ≠ ' ' Alors
            mot ← Concaténer(mot, car)
        Sinon
            Si mot ≠ "" Alors
                Ecrire(mot)
                mot ← ""
            FinSi
        FinSi
    FinPour
    Si mot ≠ "" Alors
        Ecrire(mot)
    FinSi
Fin
\end{algo}

\corrige{de l'exercice 26 (Sujet Bac 2023 C)}
\begin{algo}
Algorithme AnalyseTableau
Variables
    T : Tableau[100] : Entier
    N, i, sommePos, produitNeg, nbNuls : Entier
Début
    Ecrire("Taille du tableau : ")
    Lire(N)
    Pour i de 1 à N faire
        Ecrire("T[", i, "] = ")
        Lire(T[i])
    FinPour
    sommePos ← 0
    produitNeg ← 1
    nbNuls ← 0
    Pour i de 1 à N faire
        Si T[i] > 0 Alors
            sommePos ← sommePos + T[i]
        Sinon
            Si T[i] < 0 Alors
                produitNeg ← produitNeg * T[i]
            Sinon
                nbNuls ← nbNuls + 1
            FinSi
        FinSi
    FinPour
    Ecrire("Somme des positifs : ", sommePos)
    Ecrire("Produit des négatifs : ", produitNeg)
    Ecrire("Nombre d'éléments nuls : ", nbNuls)
Fin
\end{algo}

\corrige{de l'exercice 27 (Sujet Bac 2022 D)}
\begin{algo}
Algorithme GestionNotes
Variables
    notes : Tableau[30] : Réel
    i, nbSup : Entier
    somme, moyenne, max, min : Réel
Début
    Pour i de 1 à 30 faire
        Ecrire("Note de l'élève ", i, " : ")
        Lire(notes[i])
    FinPour
    somme ← 0
    max ← notes[1]
    min ← notes[1]
    nbSup ← 0
    Pour i de 1 à 30 faire
        somme ← somme + notes[i]
        Si notes[i] > max Alors
            max ← notes[i]
        FinSi
        Si notes[i] < min Alors
            min ← notes[i]
        FinSi
        Si notes[i] ≥ 10 Alors
            nbSup ← nbSup + 1
        FinSi
    FinPour
    moyenne ← somme / 30
    Ecrire("Moyenne : ", moyenne)
    Ecrire("Nombre d'élèves ayant ≥ 10 : ", nbSup)
    Ecrire("Note maximale : ", max)
    Ecrire("Note minimale : ", min)
Fin
\end{algo}

\corrige{de l'exercice 28 (Sujet Bac 2021 E)}
\begin{algo}
Algorithme TriEtRecherche
Variables
    T : Tableau[100] : Entier
    N, i, j, min, temp, V, debut, fin, milieu : Entier
    trouve : Booléen
Début
    Ecrire("Taille du tableau : ")
    Lire(N)
    Pour i de 1 à N faire
        Ecrire("T[", i, "] = ")
        Lire(T[i])
    FinPour
    {Tri par sélection}
    Pour i de 1 à N-1 faire
        min ← i
        Pour j de i+1 à N faire
            Si T[j] < T[min] Alors
                min ← j
            FinSi
        FinPour
        temp ← T[i]
        T[i] ← T[min]
        T[min] ← temp
    FinPour
    Ecrire("Tableau trié :")
    Pour i de 1 à N faire
        Ecrire(T[i])
    FinPour
    Ecrire("Valeur à rechercher : ")
    Lire(V)
    {Recherche dichotomique}
    debut ← 1
    fin ← N
    trouve ← Faux
    TantQue debut ≤ fin et non trouve faire
        milieu ← (debut + fin) div 2
        Si T[milieu] = V Alors
            trouve ← Vrai
        Sinon
            Si T[milieu] > V Alors
                fin ← milieu - 1
            Sinon
                debut ← milieu + 1
            FinSi
        FinSi
    FinTantQue
    Si trouve Alors
        Ecrire("Valeur trouvée à la position ", milieu)
    Sinon
        Ecrire("Valeur non trouvée")
    FinSi
Fin
\end{algo}